\chapter{Running the Calculation}
\label{ch:calculation}

\section{The Iterative Solver}

BioDYM uses an iterative convergence solver. In each pass, it processes all flows in the system:

\begin{enumerate}
  \item Check if a flow's source process has calculated inflows available
  \item If yes, apply the process logic (Splitter, Transformer, DSM, FOMP)
  \item Compute outflow values using transfer coefficients
  \item Repeat until no flow values change between passes
\end{enumerate}

This approach handles circular flows (recycling loops) naturally --- each pass propagates values further through the system until convergence.

\begin{lstlisting}[style=python]
mfa_system_results = run_mfa_calculation(
    mfa_system,
    dsm_params_config=dsm_params,      # Optional: DSM parameters
    fomp_params_config=fomp_params,     # Optional: FOMP parameters
    process_logic_map=process_logic_map,
    flow_tc_map=flow_tc_map,
    initial_stock_configs=initial_stock_configs,
)
\end{lstlisting}

\section{Process Logic Details}

\subsection{Splitter}

Distributes the total inflow to outflows using transfer coefficients. All elements maintain their input composition:

\begin{equation}
  F_{\text{out},e} = F_{\text{in},e} \times \text{TC}_{\text{material}}
\end{equation}

where $e$ is the element index and TC is the material-level transfer coefficient.

\subsection{Transformer}

Like Splitter, but allows element-specific TCs. The composition of the outflow can differ from the inflow:

\begin{equation}
  F_{\text{out},e} = F_{\text{in},e} \times \text{TC}_{e}
\end{equation}

where $\text{TC}_e$ can be different for each element.

\subsection{Pass-through}

Passes all inflow directly to a single outflow without modification:

\begin{equation}
  F_{\text{out}} = F_{\text{in}}
\end{equation}

\section{Final Balances}

After convergence, the solver calculates the stock change for every process:

\begin{equation}
  \Delta S_p = \sum F_{\text{in},p} - \sum F_{\text{out},p}
\end{equation}

For DSM processes, the stock $S$ is calculated by the DSM engine (cohort tracking), and the solver preserves it. For all other processes, $\Delta S$ is the algebraic residual.

\section{Mass Balance Verification}

After calculation, a mass balance check is displayed for each process and element:

\begin{equation}
  \text{Balance}_p = \sum F_{\text{in},p} - \sum F_{\text{out},p} - \Delta S_p
\end{equation}

This should be zero (or very close to zero, within floating-point tolerance $\sim 10^{-10}$) for every process.

\begin{warnbox}
If you see mass balance errors larger than $10^{-3}$~Mg, check your transfer coefficients and flow definitions. Common causes:
\begin{itemize}
  \item TCs not summing to 1.0 for a process
  \item Missing outflows for a process with inflows
  \item Incorrect element composition percentages
\end{itemize}
\end{warnbox}

\section{Element Hierarchy Recalculation}

After TC application, hierarchical elements are recalculated to maintain consistency. For example, if CC is defined as a percentage of DM, the solver ensures:

\begin{equation}
  \text{CC}_{\text{out}} = \text{DM}_{\text{out}} \times \frac{\text{CC}_{\text{in}}}{\text{DM}_{\text{in}}}
\end{equation}

This prevents ``transmutation'' of elements --- each flow maintains physically consistent composition.
